% TOP_PROBLEM: {"id": 1600011, "title": "Green's conjecture", "category_id": 4, "category": {"id": 4, "name": "algebra", "display_name": "Algebra", "description": "Group theory, ring theory, field theory, and algebraic structures.", "slug": "algebra", "order_index": 4, "created_at": "2026-07-31T15:26:25.671Z"}, "status": "partially_solved"}
% FIRST_POSTED: null
% ENTRY_KIND: statement_only
% Imported from: batch06.tex; UnsolvedMath ID 1600011
% Compile twice: lualatex 1600011.tex
\documentclass[11pt]{article}
\usepackage{fontspec}
\setmainfont{Latin Modern Roman}
\usepackage{amsmath,amssymb,mathtools,unicode-math}
\setmathfont{Latin Modern Math}
\usepackage{xurl,hyperref}
\hypersetup{colorlinks=true,urlcolor=blue,pdftitle={UnsolvedMath Problem 1600011}}
\newcommand{\sref}[2]{\hyperlink{source:#1:#2}{[S#2]}}
\newcommand{\cataloguescope}[1]{\paragraph{Mathematical question.}#1}
\begin{document}
% UnsolvedMath ID 1600011; source catalogue code AMR-015-0011
\setcounter{section}{1600010}
\section{Green's conjecture}\label{1600011}
{\small\sffamily UnsolvedMath ID 1600011 · Algebra}
\subsection{Definitions and mathematical statement}

\paragraph{Shared notation.}$\mathbb N=\{1,2,\ldots\}$, and $\mathbb Z,\mathbb Q,\mathbb R,\mathbb C$ denote the integers, rationals, reals and complex numbers. Any other conventions are specified locally.


Let $C$ be a smooth connected projective non-hyperelliptic curve over $\mathbb C$ of genus $g\geq3$, and let $K_C$ be its canonical line bundle. Define
\[\operatorname{Cliff}(C)=\min_{\substack{L\in\operatorname{Pic}(C)\\h^0(C,L)\geq2,\ h^1(C,L)\geq2}}\bigl(\deg L-2(h^0(C,L)-1)\bigr).
\]
For genus three use the standard value $\operatorname{Cliff}(C)=1$ when the displayed minimization set is empty. Put $V=H^0(C,K_C)$ and let $K_{p,1}(C,K_C)$ be the middle cohomology of the Koszul complex
\[\bigwedge^{p+1}V\otimes H^0(C,\mathcal O_C)\longrightarrow\bigwedge^pV\otimes H^0(C,K_C)\longrightarrow\bigwedge^{p-1}V\otimes H^0(C,K_C^{\otimes2}),
\]
whose differential alternately removes a vector and multiplies by its section; a negative exterior power is zero.
\paragraph{Statement recorded in the supplied source.}
Does every such curve satisfy
\[K_{p,1}(C,K_C)=0\quad\Longleftrightarrow\quad p\geq g-\operatorname{Cliff}(C)-1,\qquad p\geq1?
\]
Thus the length of the linear syzygy strand determines the Clifford index. \sref{1600011}{2}
\subsection{Short English statement}
Does a curve’s Clifford index coincide with the first nonlinear syzygy of its canonical embedding?
\subsection{Sources}
\begin{description}
\item[\hypertarget{source:1600011:1}{S1}] ULAM.AI and the UnsolvedMath Contributors, UnsolvedMath: A Curated Collection of Open Mathematics Problems, record 1600011 (AMR-015-0011). CC BY 4.0. \url{https://huggingface.co/datasets/ulamai/UnsolvedMath}.
\item[\hypertarget{source:1600011:2}{S2}] G. Farkas, \emph{Progress on syzygies of algebraic curves}, arXiv:1703.08056v2 (2017), §2 and §4, Definition 4.2 and Conjecture 4.4. \url{https://arxiv.org/abs/1703.08056}.
\end{description}
\label{end:1600011}
\end{document}
